\section{Materials and methods}\label{sec:methods}
Figure~\ref{fig:pipeline} summarises the design.
\begin{figure}[t]
\centering
\begin{tikzpicture}[
  font=\scriptsize, node distance=5mm and 5mm,
  box/.style={draw, rounded corners=2pt, align=center, inner sep=3pt, minimum height=9mm},
  fixed/.style={box, fill=gray!12},
  vary/.style={box, fill=blue!8, draw=blue!50!black},
  bbb/.style={box, fill=orange!12, draw=orange!60!black},
  arr/.style={-{Latex[length=2mm]}, thick}]
% row 1: shared preamble
\node[fixed] (raw) {Raw EEG\\88 subjects, 19 ch, 500\,Hz};
\node[fixed, right=8mm of raw] (pre) {Fixed preamble: 0.5--45\,Hz, 50\,Hz notch,\\avg ref, ICA + ICLabel, 128\,Hz};
\node[fixed, right=8mm of pre] (win) {10\,s windows\\(7{,}013)};
% row 2: the two varied factors, the PSWE branch and the fixed model
\node[bbb, below=11mm of raw] (pswe) {PSWE detector\\(BBB-associated)\\MPF $<6$\,Hz $\ge5$\,s};
\node[vary, right=6mm of pswe] (A) {\textbf{Factor A}: 7 node-feature sets\\P1 spectral, P2 wPLI, P3 DWT,\\P4 time, P5 raw, P6 fused, P7 PSWE};
\node[vary, right=6mm of A] (B) {\textbf{Factor B}: 3 graphs\\spatial, functional,\\hybrid};
\node[fixed, right=6mm of B] (gcn) {\textbf{Fixed} 2-layer GCN\\same hyperparameters};
% row 3: evaluation
\node[fixed, below=8mm of gcn.south east, anchor=north east] (cv) {Frozen subject-level 5$\times$5 CV (25 paired splits)\\subject = mean window logit $\to$ macro-F1};
\draw[arr] (raw) -- (pre);
\draw[arr] (pre) -- (win);
\draw[arr] (pre.south) -- ++(0,-4mm) -| (pswe.north);
\draw[arr] (win.south) -- ++(0,-4mm) -| (A.north);
\draw[arr] (win.south) -- ++(0,-4mm) -| (B.north);
\draw[arr, dashed] (pswe) -- (A);
\draw[arr] (A) -- (B);
\draw[arr] (B) -- (gcn);
\draw[arr] (gcn.south) -- (gcn.south |- cv.north);
\end{tikzpicture}
\caption{Study design. Grey components are held fixed; only the node features (Factor A) and the
graph construction (Factor B) vary, in a fully crossed $7\times3$ design evaluated on the same 25
frozen subject-level splits. The PSWE detector (orange) supplies the blood--brain-barrier-associated
marker, used as P7 node features and, in an ablation, as a co-occurrence graph.}
\label{fig:pipeline}
\end{figure}


\subsection{Dataset}
We used OpenNeuro ds004504 \cite{F1,ds004504}: resting-state, eyes-closed scalp EEG from 88
participants recorded at AHEPA University Hospital (36 AD, 23 FTD, 29 cognitively normal controls,
CN). All values below were read from the dataset's own BIDS files rather than from the descriptor
paper. Recordings use 19 electrodes of the 10--20 system (Fp1, Fp2, F3, F4, C3, C4, P3, P4, O1, O2,
F7, F8, T3, T4, T5, T6, Fz, Cz, Pz) referenced to A1--A2, sampled at \SI{500}{\hertz} on a Nihon
Kohden EEG~2100 system, with a 0.4--\SI{50}{\hertz} acquisition filter. Recording length per subject
ranges from \SI{307}{} to \SI{1291}{\second} (mean \SI{802}{\second}; \SI{19.6}{\hour} in total).
Group demographics are given in Table~\ref{tab:demographics}. The data are released under CC0; the
study uses only de-identified public data. The dataset's own metadata cite a version-1.0.9 DOI that
did not resolve when checked; we cite version 1.0.8, whose EEG data are identical (the 1.0.9
changelog records only a README edit).

\begin{table}[t]
\centering\small
\caption{Participant characteristics (mean $\pm$ SD), from \texttt{participants.tsv}. MMSE is
reported for description only and is never a model input.}
\label{tab:demographics}
\begin{tabular}{lccc}
\toprule
 & AD ($n=36$) & FTD ($n=23$) & CN ($n=29$) \\
\midrule
Age (years) & $66.4\pm7.9$ & $63.7\pm8.2$ & $67.9\pm5.4$ \\
Sex (F/M) & 24/12 & 9/14 & 11/18 \\
MMSE & $17.8\pm4.5$ & $22.2\pm2.6$ & $30.0\pm0.0$ \\
\bottomrule
\end{tabular}
\end{table}

\subsection{Fixed preprocessing}\label{sec:preproc}
Every condition shares one code path, implemented with MNE-Python \cite{mne}. Raw recordings were
band-pass filtered at 0.5--\SI{45}{\hertz}, notch filtered at \SI{50}{\hertz}, and re-referenced to
the common average. Independent component analysis (extended Infomax \cite{infomax}, 18 components
--- one fewer than the channel count because of the average reference) was fitted on a
\SI{1}{\hertz} high-passed copy of each recording, and components that ICLabel \cite{iclabel}
assigned to eye or muscle artefact with probability $\ge 0.8$ were removed. Data were then resampled
to \SI{128}{\hertz} and cut into non-overlapping \SI{10}{\second} windows (1,280 samples), giving
7,013 windows in total (30--129 per subject). ICA is fitted per recording, without labels or other
subjects' data, and therefore cannot leak information across cross-validation folds; every transform
that is estimated across subjects (feature scaling, demographic residualisation) is fitted on
training subjects only.

\subsection{Factor A: node-feature representations}\label{sec:reps}
Each window is a graph whose 19 nodes are the electrodes. Seven representations define the node
features $X\in\mathbb{R}^{19\times F}$:
\begin{description}
  \item[P1, spectral ($F=7$).] Relative power in five bands (delta 0.5--4, theta 4--8, alpha 8--13,
        beta 13--30, gamma 30--45~Hz) from Welch's method \cite{welch}
        (\SI{2}{\second} Hann segments, 50\% overlap), normalised spectral entropy, and peak
        frequency within 4--15~Hz.
  \item[P2, connectivity ($F=5$).] Node strength of the weighted phase-lag index (wPLI)
        \cite{wpli,F11} in each band: the mean wPLI between a channel and the other 18.
  \item[P3, wavelet ($F=20$).] Five-level Daubechies-4 discrete wavelet transform \cite{pywavelets};
        for the sub-bands A5, D5, D4, D3 and D2 (D1, 32--\SI{64}{\hertz}, lies above the low-pass
        cut-off), relative energy, normalised coefficient entropy, mean absolute coefficient and
        standard deviation.
  \item[P4, time domain ($F=10$).] Mean, variance, root mean square, skewness, kurtosis, zero-crossing
        rate, Hjorth mobility and complexity \cite{hjorth}, sample entropy \cite{sampen} ($m=2$,
        $r=0.2\,\sigma$) and line length. Line length replaces Hjorth activity, which equals the
        variance.
  \item[P5, raw ($F=1280$).] The window z-scored per channel.
  \item[P6, fused ($F=34$).] P1, P3 and P2 concatenated, plus each node's weighted clustering
        coefficient \cite{onnela} and eigenvector centrality in the alpha-band wPLI graph.
  \item[P7, PSWE ($F=4$).] From the PSWE detector (Section~\ref{sec:pswe}): fraction of the window
        spent in a PSWE, mean and minimum median power frequency in the window, and the local PSWE
        onset rate within $\pm\SI{60}{\second}$.
\end{description}

The wPLI between channels $i$ and $j$ was computed within each window from band-passed (4th-order
Butterworth, zero-phase) analytic signals $z_i(t)$ as
\begin{equation}
  \mathrm{wPLI}_{ij} = \frac{\bigl|\langle \operatorname{Im}(z_i z_j^{*})\rangle_t\bigr|}
                            {\langle |\operatorname{Im}(z_i z_j^{*})| \rangle_t},
\end{equation}
which is insensitive to zero-lag coupling from volume conduction \cite{wpli}.

\subsection{Factor B: graph construction}\label{sec:graphs}
Three edge definitions were crossed with every representation. \emph{Spatial}: a Gaussian kernel
$w_{ij}=\exp(-d_{ij}^2/2\sigma^2)$ on Euclidean distances between standard 10--20 electrode positions,
with $\sigma$ the median pairwise distance; this graph is identical for every window and subject.
\emph{Functional}: the alpha-band wPLI of the window. \emph{Hybrid}: the element-wise product of the
two. Each graph was sparsified by keeping the $k=4$ strongest edges per node, symmetrised as
$A=\max(A,A^\top)$, kept weighted, and normalised as $\hat A = \tilde D^{-1/2}(A+I)\tilde D^{-1/2}$.

\subsection{The fixed model}\label{sec:model}
The model is the control variable, not the contribution, so a single small graph convolutional
network (GCN) \cite{F13} was used everywhere:
$H_1=\mathrm{Dropout}_{0.5}(\mathrm{BN}(\mathrm{ReLU}(\hat A X W_0)))$,
$H_2=\mathrm{ReLU}(\hat A H_1 W_1)$, mean pooling over the 19 nodes, and a linear layer to three
classes, with $W_0\in\mathbb{R}^{F\times64}$ and $W_1\in\mathbb{R}^{64\times32}$. This gives 2,819
parameters for $F=7$ and 3,651 for $F=20$; P5 is the unavoidable exception (84,355 parameters,
because $F=1280$). Two hops of propagation already reach most of a sparsified 19-node graph, and
deeper GCNs risk over-smoothing; depth is tested as an ablation. Graph attention networks
\cite{gat} learn edge weights and would partly undo the edge definitions being compared, so a GAT
is used only as a sensitivity check.

Training used AdamW (learning rate $10^{-3}$, weight decay $10^{-4}$), batch size 64, and at most
200 epochs with early stopping (patience 20) on validation macro-F1 computed at subject level. The
loss was weighted by inverse class frequency (over subjects) divided by each subject's window count,
so every subject contributes equally regardless of recording length. A subject's prediction is the
mean of its window logits. Features were standardised with statistics from training windows only.
All hyperparameters were fixed before any result was seen and are identical in every condition.

\subsection{Evaluation protocol}\label{sec:protocol}
We used subject-level, stratified, repeated 5-fold cross-validation with five repeats (seeds 0--4),
for 25 test splits of 17--18 subjects each. Within each training set, a stratified 20\% of subjects
formed the early-stopping validation set; it was never used to tune hyperparameters. The fold
assignment was generated once, verified by assertions that training, validation and test subjects
are pairwise disjoint and that every subject is tested exactly once per repeat, and committed to the
code repository before any training run. All conditions use the same 25 splits, so every comparison
is paired. The full factorial comprises $7\times3$ cells $\times$ 25 splits $=525$ training runs.

\paragraph{Confirmation folds.} Choosing the best cell, the cells to ablate and the cell to
permute on the same splits that then estimate their performance biases those estimates upwards
(the winner's curse). We therefore generated a second, independent fold assignment with the same
procedure (seeds 5--9; again 25 splits) and committed it to the repository after the main grid had
been analysed but before any run used it. The whole grid was re-run on these folds
(\nConfirmGridRuns{} runs). The cell ranked first on seeds 0--4 is reported with its score on seeds
5--9 as a selection-free estimate, and all ablations, the leakage experiment and the permutation
test reported in the main text use the confirmation folds. Their selection-fold counterparts are
given in the Supplementary material.

The primary metric is subject-level macro-averaged F1. We also report balanced accuracy, one-vs-rest
macro ROC-AUC, PR-AUC, per-class sensitivity and specificity, Cohen's $\kappa$, the Brier score,
expected calibration error and confusion matrices. Plain accuracy is reported but never used for
ranking.

\subsection{Baselines, ablations and controls}\label{sec:ablations}
\paragraph{Baselines.} A logistic regression on age and sex alone (the demographic confound floor);
a stratified random classifier (chance); and logistic regression, RBF-kernel SVM and random forest
trained on the flattened node features of each representation with the same folds, weighting and
subject aggregation as the GCN (a non-graph twin of every representation). The MMSE is never used as
an input: control participants all score 30, so it is a proxy for the label.

The GCN learns its weights from the training subjects only and uses the validation subjects only to
choose when to stop, whereas the classical models have no stopping rule. To compare models trained
on the same subjects, we evaluate both under two regimes on the same 25 test splits:
(i)~\emph{train only}, with the classical models fitted on the training subjects and the GCN as in the
main grid; and (ii)~\emph{train + val}, with the classical models fitted on training and validation
subjects and the GCN refitted on the same subjects for the number of epochs selected by early
stopping in the corresponding main-grid run (\nRefitRuns{} refits, no validation monitoring).

\paragraph{Ablations.} Run on the two best cells of the main grid, chosen automatically from the
seed 0--4 results and evaluated on the 25 confirmation splits, each paired with the same cell's
confirmation-grid run: no graph (identity adjacency, i.e.\ a per-node MLP);
degree-preserving random graphs; $k\in\{3,6\}$ and a fully connected graph; binarised edges; depth
1 and 3; GAT instead of GCN; no subject-equal weighting; demographic residualisation (age and sex
regressed out of every feature, coefficients fitted on training windows); adding P7 (PSWE) features
to the best representation; and the PSWE co-occurrence graph (Jaccard overlap of channels' PSWE
time) as the edge type.

\paragraph{Leakage.} The best cells were re-run with \emph{window-level} stratified splitting ---
windows assigned to folds at random, ignoring subject identity --- with everything else unchanged.

\paragraph{Permutation test.} For the best cell, subject labels were permuted (all windows of a
subject keep one label) and the first repeat's five folds re-run for each permutation; the
$p$-value is $(1+\#\{\text{null}\ge\text{observed}\})/(1+N)$. The primary test uses
\cpermPThreespatialN{} permutations on the first confirmation repeat (seed 5); the original
\permPThreespatialN{} permutations on the first selection repeat (seed 0) are reported for
completeness.

\subsection{PSWE detection}\label{sec:pswe}
Following Milikovsky et al.\ \cite{B1}, with parameters as reported in \cite{B3}, we computed the
median power frequency (MPF) of each channel from FFTs of \SI{2}{\second} Hann-windowed segments with
\SI{1}{\second} overlap (one MPF value per second) over 1--\SI{45}{\hertz}. The frequency range is not
reported in the source and was set to our pass-band. A PSWE is a run of at least five consecutive
seconds with MPF $<\SI{6}{\hertz}$. Detection is per channel and per recording, and uses no labels.
For each subject we computed the PSWE rate (events per minute, averaged over channels) and the
fraction of time spent in PSWEs.

\paragraph{Background-relative PSWE (exploratory).} After observing that the fixed-threshold rate
was strongly correlated with background slowing (Section~\ref{sec:res-pswe}), we added a
post-hoc analysis, labelled as such throughout. Its definitions were fixed in the configuration
before it was run on any real data. A relative PSWE is a run of at least \SI{5}{\second} in which a
channel's MPF lies below that channel's own recording-median MPF by at least \SI{2}{\hertz}
(primary), \SI{1.5}{\hertz} or \SI{3}{\hertz}, or by at least 2 or 3 robust standard deviations
($1.4826\times$ median absolute deviation).

\subsection{Statistical analysis}\label{sec:stats}
The 25 test splits are the repeated-measures units. Main effects of representation and graph
construction, and their interaction, were tested with a balanced two-way repeated-measures ANOVA on
macro-F1, reporting partial $\eta^2$ and Greenhouse--Geisser-corrected $p$-values. Because
cross-validation scores share training data and are therefore not independent, every pairwise
comparison is reported with both the Wilcoxon signed-rank test and the Nadeau--Bengio corrected
resampled $t$-test \cite{B7}, with Holm correction \cite{holm} within
each family, Cohen's $d$ and the matched-pairs rank-biserial correlation. All 21 cells were also
compared with a Friedman test and a Nemenyi critical-difference analysis \cite{demsar}. Cell-level
95\% confidence intervals come from a subject-level bootstrap (2,000 resamples of subjects within each
repeat) and from a $t$-interval over the 25 splits. The Nadeau--Bengio correction uses the ratio of
test to training subjects of the models being compared: $17.6/56.3$ for models fitted on the training
subjects only (the main grid and the ablations) and $17.6/70.4$ when validation subjects are added
(the train + val regime).

For PSWEs, group differences were tested with Kruskal--Wallis and Dunn post-hoc tests (Holm-corrected,
with rank-biserial effect sizes), and with negative-binomial regression of event counts using
recording length as exposure, adjusted for age and sex and, in further models, for relative
delta+theta power and background MPF. The relation to slowing was assessed with Spearman and
partial (age- and sex-adjusted) rank correlations, and the relation to cognition with Spearman
correlation against MMSE within patients only. All analyses were run with SciPy, statsmodels and
custom code that is tested against reference implementations.

\subsection{Compute}
All experiments ran on free Kaggle CPU sessions (4 vCPUs); no GPU was used. The 525-run main grid
took \cpuHoursGrid{} CPU-hours, and the \nConfirmRuns{} further runs (confirmation grid, ablations and permutations on the
confirmation folds, and the train + val refits) took \cpuHoursConfirm{} CPU-hours. Every run is logged with its configuration hash, code version and seed.
